\documentclass[11pt]{article}

% ── Page geometry ──────────────────────────────────────────────────────────────
\usepackage[a4paper, top=25mm, bottom=28mm, left=20mm, right=20mm]{geometry}

% ── Fonts ──────────────────────────────────────────────────────────────────────
\usepackage[T1]{fontenc}
\usepackage{mathpazo}          % Palatino for body text and math
\usepackage{microtype}

% ── Math / science ─────────────────────────────────────────────────────────────
\usepackage{amsmath}
\usepackage{amssymb}
\usepackage{siunitx}

% ── Layout ─────────────────────────────────────────────────────────────────────
\usepackage{parskip}
\usepackage{enumitem}
\usepackage{setspace}

% ── Color (loaded before hyperref and colortbl) ────────────────────────────────
\usepackage[dvipsnames]{xcolor}
\definecolor{accent}{HTML}{1B4F8A}       % deep navy blue
\definecolor{accentlight}{HTML}{EBF2FA}  % pale blue for table header rows
\definecolor{codebg}{HTML}{F4F6F9}       % near-white blue-tinted code background
\definecolor{codenums}{HTML}{A0B4C8}     % muted blue-gray for line numbers

% ── Figures / tables ───────────────────────────────────────────────────────────
\usepackage{graphicx}
\usepackage{float}
\usepackage{booktabs}
\usepackage{array}
\usepackage{colortbl}
\usepackage[labelfont=bf, font=small]{caption}

% ── Code – tcolorbox with listings ────────────────────────────────────────────
\usepackage{listings}
\usepackage[breakable,skins,listings]{tcolorbox}

% C++ code style (used inside cppcode environment)
\lstdefinestyle{cppinner}{
  language=C++,
  basicstyle=\ttfamily\footnotesize,
  keywordstyle=\color{accent}\bfseries,
  commentstyle=\color{gray!70}\itshape,
  stringstyle=\color{BurntOrange},
  numbers=left,
  numberstyle=\tiny\color{codenums},
  numbersep=10pt,
  tabsize=4,
  breaklines=true,
  showstringspaces=false,
  keepspaces=true
}

% Plain-text format style (used inside fmtcode environment)
\lstdefinestyle{fmtinner}{
  language={},
  basicstyle=\ttfamily\small,
  numbers=none,
  tabsize=4,
  breaklines=true,
  showstringspaces=false
}

% Styled C++ code block: flat bg, accent left rule, no frame border
\newtcblisting{cppcode}[1][]{
  enhanced, breakable,
  arc=0pt, outer arc=0pt,
  colback=codebg,
  colframe=codebg,
  boxrule=0pt,
  borderline west={3pt}{0pt}{accent},
  top=5pt, bottom=5pt, left=4pt, right=4pt,
  listing only,
  listing options={style=cppinner},
  #1
}

% Styled plain-text format block: neutral bg, gray left rule
\newtcblisting{fmtcode}[1][]{
  enhanced, breakable,
  arc=0pt, outer arc=0pt,
  colback=gray!5,
  colframe=gray!5,
  boxrule=0pt,
  borderline west={3pt}{0pt}{gray!45},
  top=5pt, bottom=5pt, left=4pt, right=4pt,
  listing only,
  listing options={style=fmtinner},
  #1
}

% ── Headers / footers ──────────────────────────────────────────────────────────
\usepackage{fancyhdr}
\pagestyle{fancy}
\fancyhf{}
\fancyhead[L]{\small\color{gray!65} Sand Dune \textperiodcentered\ BOS System Documentation}
\fancyhead[R]{\small\color{gray!65} \leftmark}
\fancyfoot[C]{\small\color{gray!65} \thepage}
\renewcommand{\headrulewidth}{0.3pt}

% ── Section formatting ─────────────────────────────────────────────────────────
\usepackage{titlesec}
\titleformat{\section}
  {\large\bfseries\color{accent}}
  {\thesection}{1em}{}
  [{\vspace{2pt}\color{accent!30}\hrule height 0.6pt\vspace{4pt}}]
\titleformat{\subsection}
  {\normalsize\bfseries}
  {\thesubsection}{1em}{}
\titleformat{\subsubsection}
  {\normalsize\itshape}
  {\thesubsubsection}{1em}{}
\titlespacing{\section}{0pt}{2.5ex}{1.2ex}
\titlespacing{\subsection}{0pt}{2ex}{0.8ex}

% ── Links ──────────────────────────────────────────────────────────────────────
\usepackage[colorlinks=true, linkcolor=accent, citecolor=accent!80!black,
            urlcolor=accent!80!black]{hyperref}
\usepackage{xurl}
\urlstyle{tt}

% ── Custom commands ────────────────────────────────────────────────────────────
% Inline code: accent-colored monospace, with line-break permission after
% common delimiters so long identifiers don't overflow the right margin.
\newcommand{\codefmt}[1]{%
  \begingroup
  \ttfamily
  \hyphenchar\font=`\-
  \emergencystretch=3em
  \hspace{0pt}#1\hspace{0pt}%
  \endgroup}
\newcommand{\code}[1]{{\color{accent!85!black}\codefmt{#1}}}
\newcommand{\file}[1]{{\color{accent!70!black}\codefmt{#1}}}
\newcommand{\cls}[1]{{\bfseries\color{accent}\codefmt{#1}}}
\newcommand{\param}[1]{{\color{accent!60!black}\codefmt{#1}}}

% ── Figure borders ─────────────────────────────────────────────────────────────
% Use \figimg instead of \includegraphics to get a tight black border.
\newcommand{\figimg}[2][]{%
  \tcbox[enhanced, arc=4pt,
         boxrule=4pt, colframe=black, colback=white,
         boxsep=0pt, left=0pt, right=0pt, top=0pt, bottom=0pt]{%
    \includegraphics[#1]{#2}}%
}

% ═══════════════════════════════════════════════════════════════════════════════
\begin{document}

\begin{titlepage}
  \centering
  \vspace*{8cm}
  {\Huge\bfseries Background Oriented Schlieren SOP\par}
  \vspace{2em}
  \hrule
  \vspace{1.5em}
  {\large Xander Linzel\\
   \href{mailto:xander.linzel@irradiant.tech}{\texttt{xander.linzel@irradiant.tech}}\par}
  \vspace{1em}
  {\normalsize \today\par}

  \vspace*{8cm}
  
\end{titlepage}

    \tableofcontents

  \vspace*{1cm}
  
% ═══════════════════════════════════════════════════════════════════════════════

\section{Measurement and Image Capture}

This section is written as a step-by-step procedure. Follow the steps in order and do not change the setup between the reference image and the flow image unless the step explicitly tells you to.

\subsection{Before You Start}

\begin{enumerate}
    \item Turn on the LED backlight.
    \item Make sure the background pattern is installed correctly, meaning flat and centered on the LED panel.
    \item Make sure the printed side of the background pattern is facing down.
    \item Reduce room light as much as possible.
    \item Make sure the setup is not vibrating or being bumped.
    \item Make sure the sample holder is in place, and sturdy.
    \item Open the AmScope software on your computer and connect it to the camera.
    \item Make sure the camera is focused on the background pattern, by using teh AmScope live feed.
    \item Check that the dots in the pattern look sharp on screen.
    \item Confirm that the optical parameters are known or already calibrated:
    \begin{itemize}
        \item sensor pixel pitch $P_{px}$
        \item background-to-sample distance $Z_d$
        \item sample-to-lens distance $Z_a$
        \item focal length $f$
        \item sample refractive index $n$
        \item sample thickness $t$
    \end{itemize}
\end{enumerate}

\subsection{Reference and Flow Image Capture}

\begin{enumerate}
    \item Make sure the live camera image is visible.
    \item Check focus again.
    \item Check brightness again.
    \item Make sure there is \textbf{no sample} in the optical path.
    \item Press \textbf{Capture Photo}.
    \item Press \textbf{Take Photo}.
    \item Save this image using the correct reference name.
    \item Use a clear file name such as \texttt{sampleA\_ref.png}.
    \item Place the sample in the sample holder.
    \item Make sure the sample is seated properly.
    \item Make sure the sample is not heavily tilted.
    \item Press \textbf{Capture Photo} again.
    \item Press \textbf{Take Photo}.
    \item Save this image using the correct flow name.
    \item Use a clear file name such as \texttt{sampleA\_flow.png}.
    \item If multiple flow images are needed for the same reference, repeat the previous capture step and save each flow image with a unique name.
    \item Do not change focus, camera position, or geometry between the reference image and its matching flow images.
\end{enumerate}

\subsection{Practical Notes}

\begin{itemize}
    \item Use one reference image and its matching flow images as one measurement group.
    \item If you change the reference condition, take a new reference image and make a new group in the software.
    \item Strong ambient light, vibration, blur, or sample tilt can create false displacement signals.
    \item If the dot pattern looks soft or smeared, stop and refocus before continuing.
\end{itemize}

\section{Using the Application}
\label{sec:usage}

\subsection{Typical Workflow}

\begin{enumerate}
    \item Open the \textbf{Sand Dune} application.
    \item Go to the \textbf{Load Images} panel.
    \item If this is the first measurement group, use the default group already shown.
    \item If you need another reference with its own matching flow images, press \textbf{Add} to create a new group.
    \item Use the group controls to move to the group you want to work on.
    \item In the \textbf{Reference Image} section, press \textbf{Browse}.
    \item Select the reference image for the active group.
    \item In the \textbf{Flow Images} section, press \textbf{Browse}.
    \item Select one or more flow images that belong to that same reference.
    \item If the active group contains more than one flow image, use the \textbf{Active Flow} controls to move between them.
    \item Repeat the group / reference / flow loading steps for every measurement set.
    \item Go to the \textbf{Parameters} panel.
    \item Enter the correlation settings:
    \begin{itemize}
        \item window size
        \item overlap
        \item PID on/off
        \item PID iterations
        \item PID relaxation
        \item PID smoothing passes
    \end{itemize}
    \item Enter the optical settings:
    \begin{itemize}
        \item sensor pixel pitch
        \item background-to-sample distance
        \item sample-to-lens distance
        \item lens focal length
        \item aperture diameter
    \end{itemize}
    \item Enter the sample settings:
    \begin{itemize}
        \item thickness
        \item refractive index
    \end{itemize}
    \item Choose the reconstruction mode:
    \begin{itemize}
        \item RI variation
        \item thickness variation
    \end{itemize}
    \item Choose the integration method:
    \begin{itemize}
        \item Masked Poisson
        \item Frankot--Chellappa
    \end{itemize}
    \item Turn \textbf{Field Correction} on or off as needed.
    \item If reconstruction should be limited to the sample region, enable the mask.
    \item Set the mask position and radius.
    \item Check the \textbf{Calculations} panel.
    \item Make sure the displayed grid spacing, field width, and optical values look reasonable.
    \item Go to the \textbf{Pipeline} panel.
    \item Press \textbf{Run Full Pipeline} to process everything in order.
    \item If needed, run the stages one at a time:
    \begin{itemize}
        \item Correlation
        \item Validation
        \item Reconstruction
    \end{itemize}
    \item Wait for the status to show \textbf{Done}.
    \item Open the visualization panels and inspect:
    \begin{itemize}
        \item Correlation
        \item Validation
        \item Surface
    \end{itemize}
    \item For correlation and validation, choose the map you want to view:
    \begin{itemize}
        \item $u$
        \item $v$
        \item magnitude
        \item s2n
    \end{itemize}
    \item If needed, toggle between raw pixel display and scaled physical units.
    \item Adjust the colormap min and max sliders if the range needs to be tightened.
    \item Go to the \textbf{Save} panel.
    \item Choose the output directory.
    \item Enter an output name if desired.
    \item Press \textbf{Save All}.
\end{enumerate}

\subsection{Panel Reference}

\subsubsection{Load Panel}

The \textbf{Load Images} panel is now group-based.

\begin{itemize}
    \item \textbf{Groups section:} Shows the active group number and whether the active group has a reference image and flow images loaded.
    \item \textbf{$<$ Group / Group $>$:} Moves between groups.
    \item \textbf{Add:} Creates a new empty group.
    \item \textbf{Delete:} Deletes the active group. If only one group remains, it is reset to an empty group.
    \item \textbf{Reference Image section:} Loads the reference image for the active group.
    \item \textbf{Flow Images section:} Loads one or more flow images for the active group.
    \item \textbf{Active Flow section:} Appears only when the active group has more than one flow image loaded.
\end{itemize}

\subsubsection{Parameters Panel}

\begin{center}
\begin{tabular}{lll}
  \toprule
  \rowcolor{accentlight} Group & Control & Effect \\
  \midrule
  Correlation & Window size & Interrogation window side length (pixels) \\
              & Overlap     & Window overlap (pixels; step = size $-$ overlap) \\
              & Window Deformation (PID) & Enables iterative window deformation \\
              & PID iterations & Number of PID passes after baseline \\
              & PID relaxation & Damping on PID residual updates \\
              & PID smoothing  & Number of smoothing passes on the PID field \\
  \midrule
  Optical     & $P_{px}$       & Sensor pixel pitch (\si{\micro\metre}) \\
              & $Z_d$          & Background-to-sample distance (\si{mm}) \\
              & $Z_a$          & Sample-to-lens distance (\si{mm}) \\
              & $f$            & Lens focal length (\si{mm}) \\
              & $d_a$          & Aperture diameter (\si{mm}) \\
  \midrule
  Sample      & $t$            & Sample thickness (\si{mm}) \\
              & $n$            & Sample refractive index \\
  \midrule
  Reconstruction & Mode        & RI variation or thickness variation \\
                 & Solver      & Masked Poisson or Frankot--Chellappa \\
                 & Field Correction & Applies the refraction correction field \\
  \midrule
  Mask        & Mask Enabled   & Enables circular ROI masking \\
              & X / Y Position & Mask centre location \\
              & Radius         & Mask radius \\
  \midrule
  Background  & Dot diameter   & Background dot diameter estimate (\si{mm}) \\
  \bottomrule
\end{tabular}
\end{center}

\subsubsection{Calculations Panel}

The \textbf{Calculations} panel updates live from the current parameters and shows derived values such as:
\begin{itemize}
    \item sensor image distance
    \item magnification
    \item effective dot size in pixels
    \item suggested correlation window size
    \item spatial resolution on the sample plane
\end{itemize}

Use this panel as a quick sanity check before running the pipeline.

\subsubsection{Pipeline Panel}

The \textbf{Pipeline} panel shows:
\begin{itemize}
    \item stage status
    \item run buttons for each stage
    \item a \textbf{Run Full Pipeline} button
    \item progress bars
    \item ETA values during processing
\end{itemize}

The pipeline processes all loaded groups and all flow images currently in the session.

\subsubsection{Visualization Panel}

\begin{itemize}
    \item \textbf{Correlation panel:} Shows $u$, $v$, magnitude, or s2n.
    \item \textbf{Validation panel:} Shows $u$, $v$, magnitude, or s2n after validation.
    \item \textbf{Surface panel:} Shows the reconstructed scalar field.
    \item \textbf{Raw/scaled toggle:} Switches correlation and validation display between raw pixels and scaled physical units.
    \item \textbf{Min/max sliders:} Control the displayed colormap range.
\end{itemize}

\subsubsection{Save Panel}

The \textbf{Save} panel writes output for all processed groups and flows.

If a name is entered in the UI, saved files are named like:
\begin{fmtcode}
<ui_name>_<flow_name>_g<group>_f<flow>_correlation.csv
<ui_name>_<flow_name>_g<group>_f<flow>_val.csv
<ui_name>_<flow_name>_g<group>_f<flow>_surface.csv
\end{fmtcode}

If no name is entered in the UI, the flow name is used:
\begin{fmtcode}
<flow_name>_g<group>_f<flow>_correlation.csv
<flow_name>_g<group>_f<flow>_val.csv
<flow_name>_g<group>_f<flow>_surface.csv
\end{fmtcode}

\subsubsection{Settings Panel}

The \textbf{Settings} panel contains the light/dark theme switch.

\subsection{CSV File Formats}

\paragraph{Correlation and Validation CSV.}
The correlation and validation result files share the following layout:
\begin{fmtcode}
rows, cols
<height>, <width>
u, v, s2n
<u00>, <v00>, <s2n00>
...
\end{fmtcode}
Values are saved from the displayed processed fields.

\paragraph{Surface CSV.}
The surface file uses column-major ordering:
\begin{fmtcode}
rows, cols
<height>, <width>
<z[0,0]>
<z[1,0]>
...
\end{fmtcode}
Values are stored in column-major order.



\end{document}
